Die Erde dreht sich in $T = \TEqO$ einmal um ihre Achse. Ein Punkt am Äquator legt dabei den Umfang $2\pi r$ zurück:
\begin{align}
 v &= \vEqF \\
   &= \frac{2\pi\cdot\rEq}{\TEq} \\
   &= \vEq \approx \result{\vEqP{0}{3}}
\end{align}
Genau genommen ist die Rotationsdauer der Erde gegenüber den Fixsternen (siderischer Tag) etwas kürzer als ein Tag; der Unterschied ist hier vernachlässigbar.
